\chapter{Mở Tiết Bắt Lớp}
\chapterintro{Đầu giờ không cần nói dài}{Một câu gọn và có duyên thường kéo lớp về nhanh hơn một đoạn nhắc trật tự.}

\section{Hôm nay ai định học tử tế?}
\begin{manthoai}{Màn thoại}
\textbf{Thầy/Cô:} Hôm nay ai định học tử tế? Giơ tay thật đi.\\
\textbf{Học sinh:} Cả lớp luôn ạ.\\
\textbf{Thầy/Cô:} Tốt. Thế mình bắt đầu bằng việc tử tế đầu tiên: nhìn lên bảng.
\end{manthoai}
\begin{dungkhinào}
Dùng lúc lớp vừa vào còn nhốn nháo nhẹ. Câu này khiến học sinh bật cười nhưng cũng tự bước vào khuôn rất nhanh.
\end{dungkhinào}
\begin{nhipthem}
Có thể đổi thành: \textit{``Ai định làm thầy cô vui trong 5 phút đầu?''}
\end{nhipthem}
\ngancach

\section{Lớp mình tỉnh chưa hay còn đang loading?}
\begin{manthoai}{Màn thoại}
\textbf{Thầy/Cô:} Lớp mình tỉnh chưa hay còn đang loading?\\
\textbf{Học sinh:} Còn lag ạ.\\
\textbf{Thầy/Cô:} Rồi, mình khởi động nhẹ để hết lag trước khi học thật.
\end{manthoai}
\begin{dungkhinào}
Dùng với lớp có thói quen vui, hợp cho đầu giờ sáng hoặc sau giờ ra chơi.
\end{dungkhinào}
\begin{nhipthem}
Nối ngay bằng một câu hỏi mở hoặc một mini game 30 giây để tránh câu nói chỉ dừng ở chỗ gây cười.
\end{nhipthem}
\ngancach

\section{Ai nói chuyện tiếp là phải nói luôn ý chính của bài}
\begin{manthoai}{Màn thoại}
\textbf{Thầy/Cô:} Ai nói chuyện tiếp là phải nói luôn ý chính của bài.\\
\textbf{Học sinh:} Dạ thôi ạ.\\
\textbf{Thầy/Cô:} Thế thì cả lớp nghe ý chính từ thầy cô trước đã.
\end{manthoai}
\begin{dungkhinào}
Dùng khi cần cắt nói chuyện riêng mà không muốn gắt.
\end{dungkhinào}
\begin{nhipthem}
Nếu lớp mạnh miệng, có thể mời thật một em nói thử rồi khen tinh thần trước khi quay lại bài.
\end{nhipthem}
\ngancach

\section{Hôm nay lớp muốn được gọi là lớp tỉnh hay lớp cần cứu?}
\begin{manthoai}{Màn thoại}
\textbf{Thầy/Cô:} Hôm nay lớp muốn được gọi là lớp tỉnh hay lớp cần cứu?\\
\textbf{Học sinh:} Lớp tỉnh ạ.\\
\textbf{Thầy/Cô:} Tốt. Danh hiệu đó giữ giúp thầy cô ít nhất 10 phút đầu nhé.
\end{manthoai}
\begin{dungkhinào}
Hợp khi đầu tiết lớp hơi lỏng mà thầy cô muốn kéo về bằng một câu vừa vui vừa có yêu cầu.
\end{dungkhinào}
\begin{nhipthem}
Có thể đổi vế sau thành: \textit{``Thế thì nhìn lên bảng để chứng minh nào.''}
\end{nhipthem}
\ngancach

\section{Ai đang tới lớp bằng thân người nhưng não còn ở ngoài sân?}
\begin{manthoai}{Màn thoại}
\textbf{Thầy/Cô:} Ai đang tới lớp bằng thân người nhưng não còn ở ngoài sân?\\
\textbf{Học sinh:} Có ạ.\\
\textbf{Thầy/Cô:} Rồi, mình kéo não vào trước rồi học tiếp.
\end{manthoai}
\begin{dungkhinào}
Rất hợp sau giờ ra chơi hoặc đầu buổi sáng khi học sinh ngồi vào chỗ nhưng chưa thật sự nhập tiết.
\end{dungkhinào}
\begin{nhipthem}
Nói xong thì chuyển ngay sang một câu hỏi cực ngắn để lớp vào bài, không để câu vui đứng một mình.
\end{nhipthem}
